\section{El recorrido, momento por momento}
El trámite avanza en cuatro momentos: 1. Reúne tus requisitos, 2. Propón sínodo y fecha, 3. Invitaciones al sínodo, 4. Examen y cierre. Es igual en todas las modalidades; las pantallas son de la modalidad \emph{Tesis}.
\subsection{Momento 1. Reúne tus requisitos}
Capturas tus datos (incluida la modalidad y los datos de tu director), pides el voto aprobatorio, firmas tu formato de datos y subes cada requisito en PDF. Apoyo a Grado los valida y, con los originales, hace el cotejo.
\subsubsection*{Tus requisitos}
\begin{multicols}{2}\small
\textbf{Identidad y antecedentes académicos}
\begin{itemize}
  \item Acta de nacimiento
  \item Certificado parcial del posgrado
  \item Certificado total de licenciatura
  \item Certificado total de maestría (solo doctorado)
  \item Cédula profesional de licenciatura
  \item Cédula profesional de maestría (solo doctorado)
  \item CURP
\end{itemize}
\textbf{Pagos}
\begin{itemize}
  \item Pago de derechos de examen recepcional
  \item Pago de expedición de título de grado
  \item Hoja de referencia de pago del certificado total
\end{itemize}
\textbf{Trabajo de grado}
\begin{itemize}
  \item Oficio-constancia de registro de protocolo de tesis
  \item Certificado de autenticidad del trabajo
  \item Evidencia de artículo aceptado o capítulo de libro (solo doctorado)
  \item Autorización para continuar con los trámites de obtención de grado (solo si aplica el art. 53 del REA)
\end{itemize}
\textbf{Constancias y no adeudos}
\begin{itemize}
  \item Constancia de no adeudo bibliotecario
  \item Oficio de donación del libro
  \item Constancia de no adeudo de mobiliario y equipo — lo genera la plataforma
  \item Constancia del llenado del CVU
  \item Folio del Sistema de Seguimiento de Egresados de posgrado
\end{itemize}
\textbf{Entregas en físico}
\begin{itemize}
  \item 3 memorias USB con tu trabajo de grado — en persona
  \item Resumen impreso del trabajo — en persona
  \item 7 fotografías tamaño óvalo miñón — en persona
  \item 2 fotografías tamaño semitítulo — en persona
\end{itemize}
\textbf{Formatos}
\begin{itemize}
  \item Voto aprobatorio de tu director — se firma en línea
  \item Formato de datos del interesado — lo genera la plataforma
  \item Formato de revisión para autorización de evaluación — en persona
  \item Formato de integración de expediente — lo genera la plataforma
\end{itemize}
\end{multicols}
\captura{tesis_egresado_05_siguiente_datos}{Tu siguiente paso: siempre una sola acción, aquí «Capturar mis datos».}
\captura{tesis_egresado_06_datos}{Tus datos: aquí eliges la modalidad («Tesis») y capturas los datos de tu director con su correo.}
\captura{tesis_egresado_07_siguiente_voto}{Con tus datos guardados, el siguiente paso cambia: pedir el voto aprobatorio a tu director.}
\captura{tesis_egresado_08_voto}{Pides el voto aprobatorio: tu director recibe por correo su enlace para firmarlo en línea.}
\captura{tesis_egresado_09_firmar_datos}{Revisas tu formato de datos del interesado y lo firmas en línea con un botón.}
\captura{tesis_egresado_10_subir}{Subes cada requisito en PDF desde su renglón; queda «En revisión».}
\captura{tesis_egresado_11_siguiente_revision}{Con todo subido, el siguiente paso es de Apoyo a Grado: revisar tus documentos.}
\captura{tesis_egresado_13_devuelto}{Si Apoyo a Grado te devuelve un documento, tu siguiente paso te lleva a corregirlo.}
\captura{tesis_egresado_14_corregir}{Documento devuelto: lees la observación de Apoyo a Grado, eliges el PDF corregido y pulsas «Reemplazar».}
\subsection{Momento 2. Propón sínodo y fecha}
Capturas a los cinco propietarios y los dos suplentes con su correo, la fecha y la hora. Al enviar, tu director firma la solicitud y cada integrante su visto bueno, desde su correo.
\captura{tesis_egresado_15_siguiente_sinodo}{Momento 2: propones a tu sínodo y la fecha de tu examen.}
\captura{tesis_egresado_16_sinodo}{Propuesta de sínodo: cinco propietarios y dos suplentes, cada uno con su correo; marca a los externos.}
\captura{tesis_egresado_17_fecha}{Fecha y hora propuestas: la plataforma solo acepta fechas con al menos 10 días hábiles.}
\captura{tesis_egresado_18_esperando_firmas}{Mientras firman, ves quién falta; puedes reenviar el enlace o corregir un correo.}
\captura{tesis_egresado_19_vigencia_aviso}{Tu siguiente paso te avisa qué sinodal debe renovar su firma electrónica UAEMéx.}
\subsection{Momento 3. Invitaciones al sínodo}
Apoyo a Grado emite las invitaciones, la Coordinación las firma y cada integrante acusa de recibido en línea.
\captura{tesis_egresado_20_acuses}{Momento 3: ves quién ya acusó de recibido su invitación.}
\subsection{Momento 4. Examen y cierre}
Presentas tu examen; aprobado, registras tu grado en el SICDE y capturas el folio. Apoyo a Grado emite las constancias e integra tu expediente.
\captura{tesis_egresado_21_sicde}{Aprobado el examen, registras tu grado en el SICDE (48 horas) y capturas aquí el folio.}
\captura{tesis_egresado_22_concluido}{Trámite concluido: los oficios y constancias quedan en tu expediente.}
\section{Desde el celular}
El trámite funciona igual en el teléfono: el siguiente paso arriba y los momentos plegados.
\capturasmovil{tesis_egresado_03_movil_modalidad}{Desde el celular: la página de la modalidad y su manual.}{tesis_egresado_12_movil_tramite}{Desde el celular: tu siguiente paso y los cuatro momentos.}
